\begin{table}[tp]
\centering
\caption{Convex optimization under generalized smoothness (\href{https://github.com/k3nfalt/longtails/blob/main/LongTails/Setups/GeneralizedSmoothAndConvex.lean}{GeneralizedSmoothAndConvex}).}\label{tab:GeneralizedSmoothAndConvex}
\parbox{\linewidth}{\footnotesize Convex $f:E\to\mathbb R$ on a real Hilbert space has a differentiable gradient with $\|\nabla^2 f(x)\|\le L_0+L_1\|\nabla f(x)\|$ everywhere. Here $L_0,R,\Delta,\varepsilon>0$, $L_1\ge0$, and a minimizer $x_\star$ satisfies $\|x_0-x_\star\|\le R$ and $f(x_0)-f(x_\star)\le\Delta$. The oracle returns exact gradients; the target is $f(x_K)-f(x_\star)\le\varepsilon$. Our constructions use $L_0,L_1,R,\Delta$, return a deterministic last iterate, and count every gradient call.}\par\smallskip
\begin{threeparttable}
\footnotesize
\setlength{\tabcolsep}{4pt}\renewcommand{\arraystretch}{1.25}
\begin{tabularx}{\linewidth}{@{}>{\raggedright\arraybackslash}p{0.20\linewidth}>{\raggedright\arraybackslash}p{0.32\linewidth}>{\raggedright\arraybackslash}p{0.25\linewidth}>{\raggedright\arraybackslash}X@{}}
\toprule
\textbf{Result} & \textbf{Oracle complexity} & \textbf{Novelty} & \textbf{Limitations} \\
\midrule
\textbf{(1)} \href{https://github.com/k3nfalt/longtails/blob/main/LongTails/Theorems/GeneralizedSelfNormalizedConvex.lean}{Generalized\allowbreak Self\allowbreak Normalized\allowbreak Convex} \textbf{(Ours)} $(\star\star\star)$ \href{https://github.com/k3nfalt/longtails/blob/main/scripts/reviews/GeneralizedSelfNormalizedConvex.md}{(review)}
& $\begin{aligned}\mathcal{O}\bigl(1+\sqrt{L_0R^2/\varepsilon}\qquad\\{}+\min\{\lceil q^{2/3}\rceil\ell,\,\ell+(q\ell)^{2/3}\}\bigr)\end{aligned}$ gradients
& & \\
\midrule
\textbf{(2)} \href{https://github.com/k3nfalt/longtails/blob/main/LongTails/Theorems/GeneralizedAcceleratedConvex.lean}{Generalized\allowbreak Accelerated\allowbreak Convex} \textbf{(Ours)} $(\star\star\star)$ \href{https://github.com/k3nfalt/longtails/blob/main/scripts/reviews/GeneralizedAcceleratedConvex.md}{(review)}
& $\mathcal{O}\bigl(1+\sqrt{L_0R^2/\varepsilon}+\lceil q^{2/3}\rceil\ell\bigr)$ gradients
& & \\
\midrule[1.2pt]
\textbf{(3)} AGD \citep{Tyurin2026v2CompositeTable}, Thm.~4.3.\par On $\mathbb R^d$.
& $\begin{aligned}\mathcal{O}\bigl(1+\sqrt{L_0R^2/\varepsilon}\qquad\\{}+q\log(1+\min\{L_1^2\Delta/L_0,\Delta/\varepsilon\})\\{}+\log(1+\Delta/(L_0R^2))\bigr)\end{aligned}$ gradients.\par Derived bound: $\bar R=R$, $\Gamma_0=2\Delta/R^2$, using a flat-coordinate extension to realize distance $R$.
& \textbf{(1)}: no larger order. For $q\to\infty$, $\Delta=L_0R^2$, $\varepsilon=L_0R^2q^{-4/3}$:\par \textbf{(3)}: $\mathcal{O}(q\log q)$;\par \textbf{(1)}: $\mathcal{O}((q\log q)^{2/3})$;\par \textbf{(2)}: $\mathcal{O}(q^{2/3}\log q)$.\par \textbf{(2)}: moderate-accuracy gain, not uniform.
& \textbf{(2)}: unsaturated logarithm can cost more when $\varepsilon<L_0/L_1^2$ and $q^{2/3}\ell>\sqrt{L_0R^2/\varepsilon}$.\par \textbf{(1)}, \textbf{(2)}: two gradients per iteration instead of one. \\
\midrule
\textbf{(4)} AGMsDR \citep{gscVankov2025v3}, Thm.~6.2.\par On $\mathbb R^d$, $f\in C^2$.
& $\begin{aligned}\mathcal{O}\bigl((\nu+1)[1+\sqrt{L_0R^2/\varepsilon}\\{}+\lceil q^{2/3}\rceil\ell]\bigr)\end{aligned}$ value/gradient calls, including exact segment minimizations; $\nu$ has no parameter bound.
& \textbf{(1)}, \textbf{(2)}: gradients only; no exact segment search. \textbf{(2)} matches iterations. \textbf{(1)} saves $\ell^{1/3}$ iterations when $q\ge\sqrt\ell$, $q^{2/3}\ell\gg1+\sqrt{L_0R^2/\varepsilon}$. Call bounds are not comparable.
& \textbf{(1)}, \textbf{(2)}: need $R,\Delta$ in addition for initialization. \\
\midrule
\textbf{(5)} Accelerated GRAAL \citep{gscBorodichKovalev2025v2}, Cor.~3.\par On $\mathbb R^d$, $f\in C^2$, $L_1>0$.
& $\begin{aligned}\mathcal{O}\bigl(1+\sqrt{L_0R^2/\varepsilon}+q^3\\{}+(1+q^2)\log(1/(\eta_0L_0))\bigr)\end{aligned}$ gradient and value calls; $0<\eta_0\le e^{-q}/L_0$.
& \textbf{(1)}, \textbf{(2)}: no function values. At $\eta_0=e^{-q}/L_0$, $\Delta=L_0R^2$, $\varepsilon=L_0R^2q^{-4}$, $q\to\infty$: total $\mathcal{O}(q^3)\to\mathcal{O}(q^2)$. \textbf{(2)} replaces cubic startup by $\lceil q^{2/3}\rceil\ell$; gain depends on accuracy.
& \textbf{(1)}, \textbf{(2)}: need $L_0,L_1,R,\Delta$; \textbf{(5)} updates need only admissible $\eta_0$. Startup in \textbf{(1)}, \textbf{(2)} depends on accuracy. \\
\bottomrule
\end{tabularx}
\begin{tablenotes}[flushleft]\footnotesize
\item[] \textbf{Abbr:} $q=L_1R$; $\ell=1+\log(1+\Delta/\varepsilon)$; $\eta_0$ is the initial step size; $\nu$ counts oracle calls per exact segment minimization. Complexities are total upper bounds, up to universal factors; value and gradient evaluations each count as primitive calls.
\end{tablenotes}
\end{threeparttable}
\end{table}
\clearpage
